
\subsubsection{5.1 H-E1: Behavioral Variance Exists}

\begin{table}[htbp]
\centering
\resizebox{\columnwidth}{!}{%
\begin{tabular}{llll}
\toprule
Metric & Value & Threshold & Status \\
\midrule
Total Variance & 0.0842 & --- & --- \\
Residual Variance & 0.0569 & --- & --- \\
Residual Ratio & **0.6758** & > 0.05 & **PASS** \\
$R^2$ of Stratified Baseline & 0.3242 & --- & --- \\
Factor Above Threshold & 13.$5\times$ & --- & --- \\
\bottomrule
\end{tabular}
}
\end{table}

67.6\% of class-wise variance is unexplained by overall accuracy.

\textbf{Per-Class Analysis:}

\begin{table}[htbp]
\centering
\resizebox{\columnwidth}{!}{%
\begin{tabular}{llll}
\toprule
Class & Name & Mean Accuracy (Difficulty) & Variance \\
\midrule
0 & airplane & 0.196 & 0.080 \\
1 & automobile & 0.465 & 0.163 \\
2 & bird & 0.080 & 0.057 \\
3 & cat & 0.029 & 0.007 \\
4 & deer & 0.039 & 0.013 \\
5 & dog & 0.182 & 0.059 \\
6 & frog & 0.138 & 0.044 \\
7 & horse & 0.212 & 0.057 \\
8 & ship & 0.217 & 0.081 \\
9 & truck & 0.309 & 0.127 \\
\bottomrule
\end{tabular}
}
\end{table}

Automobile (variance 0.163) and truck (variance 0.127) show the highest per-class variance, indicating models differentiate substantially on vehicle recognition capability.

\textbf{Model Statistics:}
\begin{itemize}
\item Mean overall accuracy: 18.7\%
\item Standard deviation of overall accuracy: 9.1\%
\end{itemize}

The low mean accuracy (18.7\%) indicates models in this subset are undertrained, which is expected for a zoo containing models across different training stages and hyperparameter configurations.

\subsubsection{5.2 H-M1: Weight Statistics Fail}

\begin{table}[htbp]
\centering
\resizebox{\columnwidth}{!}{%
\begin{tabular}{llll}
\toprule
Metric & Value & Threshold & Status \\
\midrule
Proposed (Weight Features) Mean $R^2$ & **-0.084** & --- & --- \\
Baseline (Stratified) Mean $R^2$ & **0.117** & --- & --- \\
$\Delta R^2$ & **-0.201** & > 0 & **FAIL** \\
\bottomrule
\end{tabular}
}
\end{table}

Simple weight statistics perform worse than the stratified baseline.

\textbf{Per-Class $R^2$ Comparison:}

\begin{table}[htbp]
\centering
\resizebox{\columnwidth}{!}{%
\begin{tabular}{llll}
\toprule
Class & Baseline $R^2$ & Proposed $R^2$ & Winner \\
\midrule
0 & -0.014 & 0.202 & Proposed \\
1 & 0.187 & 0.220 & Proposed \\
2 & -0.111 & 0.112 & Proposed \\
3 & -0.025 & -0.041 & Baseline \\
4 & -0.090 & **-2.739** & Baseline \\
5 & 0.276 & 0.252 & Baseline \\
6 & 0.187 & 0.231 & Proposed \\
7 & 0.286 & 0.158 & Baseline \\
8 & 0.135 & 0.262 & Proposed \\
9 & 0.336 & 0.501 & Proposed \\
\bottomrule
\end{tabular}
}
\end{table}

Class 4 (deer) shows extreme negative $R^2$ = -2.74 for the proposed model, indicating severe overfitting or feature-target misalignment for this class. The proposed method outperforms baseline on 6 of 10 classes individually, but the large failure on class 4 dominates the mean.
